\documentclass[10pt,a4paper]{amsart}
\usepackage[margin=19mm]{geometry}
\usepackage{amsmath,amssymb,mathtools}
\usepackage[scr=rsfs,cal=euler]{mathalfa}
\usepackage{xfrac}
\usepackage[unicode,hidelinks,bookmarks=false]{hyperref}
\newcommand{\ie}{i.e.\ }
\newcommand{\cf}{cf.\ }
\newcommand{\resp}{respectively\ }
\newcommand{\Subsection}{\subsection}
\providecommand{\nobreakdash}{\nobreak\hbox{-}}
\setlength{\emergencystretch}{2em}
\multlinegap0pt
\usepackage{amssymb}
\usepackage{mathtools}
\usepackage[shortlabels]{enumitem}
\newtheorem{theorem}{Theorem}
\newtheorem{proposition}{Proposition}
\newcommand{\E}{\mathbb{E}}
\newcommand{\R}{\mathbb{R}}
\DeclareMathOperator{\var}{Var}
\title{Limit theorems for $p$-domain functionals of stationary Gaussian fields}
\author{Nikolai Leonenko, Leonardo Maini, Ivan Nourdin, Francesca Pistolato}
\date{Source: arXiv:2402.16701v2 (CC BY 4.0)}
\begin{document}
\maketitle

\noindent\emph{Source and licence.} Limit theorems for $p$-domain functionals of stationary Gaussian fields. Version arXiv:2402.16701v2 (CC BY 4.0), distributed by arXiv under the Creative Commons Attribution 4.0 International licence, http://creativecommons.org/licenses/by/4.0/, as recorded in the arXiv licence metadata for this exact version and shown on the abstract page https://arxiv.org/abs/2402.16701v2. Full licence notice: \texttt{licenses/arxiv-2402.16701-CC-BY-4.0.txt}.\par\smallskip

\noindent\textit{Editorial scope.} Counted unit: the article's proposition prop:quantHq (source0.tex lines 565--573). The article's complete original proof of it is delivered with it (source0.tex lines 575--607, one \texttt{proof} environment of 33 lines). No mathematics is written, repaired or re-derived: the statement and the proof are the article's own lines, in the article's own notation, and no retained line is substituted or re-flowed.  Retained uncounted as the source-local context the proof needs: lines 147--150; lines 371--385; lines 458--470; lines 459--461; lines 478--482; lines 483--485; lines 486--488.  Nothing was omitted from any retained span, so the package carries no omission record.  No retained line contains a citation, so the wrapper carries no bibliography and no bibliography entry.  No retained line contains a figure environment, an image-inclusion command or drawing code, so no image is delivered and the package declares no asset dependency.  Editorial formatting only, in addition: an AMS \texttt{amsart} preamble that restates the article's own declarations and macros used by the retained lines, namely the environments \texttt{theorem}, \texttt{proposition} on separate counters, together with the standard packages those lines need.  The retained environments print in this document as Theorem 1, Proposition 1 in order of appearance, whereas the article prints the same items with the numbering of its own full document.  The complete native source of the article as distributed in the arXiv e-print for this exact version is bundled unchanged as \texttt{sources/155154/source0.tex}.
\begin{equation}
    \label{equ:sep}
    C(x_1,\ldots,x_p)= \prod_{i=1}^p C_i(x_i),\quad x_i\in \R^{d_i},\quad 1\leq i\leq p,
\end{equation}

\begin{theorem}\label{thm:4thQuant}
Fix $q\ge 2$ and consider $h\in \mathcal H^{\odot q}$. 
Then
\begin{eqnarray}
\E[I_q(h)^4]-3\E[I_q(h)^2]^2
&=&\frac3q\sum_{r=1}^{q-1} rr!^2\binom{q}{r}^4(2q-2r)!
\|h\widetilde{\otimes}_r  h\|^2_{2q-2r}\\
&=&\sum_{r=1}^{q-1} q!^2\binom{q}{r}^2\left\{\|h\otimes_r h\|^2_{2q-2r}+\binom{2q-2r}{q-r}\|h\widetilde{\otimes}_r h\|^2_{2q-2r}\right\},
\end{eqnarray}
and, with $N\sim N(0,1)$,
\begin{align}
    &d_{TV}(I_q(h),N)\le \frac{2}{\sqrt{3}}\sqrt{\E[I_q(h)^4]-3\E[I_q(h)^2]^2},\\
    &{d_{W}(I_q(h),N)\le \frac{2}{\sqrt{3\pi}}\sqrt{\E[I_q(h)^4]-3\E[I_q(h)^2]^2}.} 
\end{align}
\end{theorem}

\begin{proposition}\label{prop:suffHq}
    Let $B=(B_x)_{x\in\R^d}$ be a real-valued, continuous, centered, stationary Gaussian field with unit-variance. Let $C:\R^d \rightarrow \R$ be the covariance function of $B$ and assume it is separable in the sense of \eqref{equ:sep}. Then \begin{equation}\label{equ:varHq}
        \var(Y(t_1,\ldots,t_p)[q]) = (q!)^{1-p} \prod_{i=1}^p \var (Y_i{(t_i)}[q]).
    \end{equation}
    Moreover, the following holds: if there exists $i\in\{1,\ldots,p\}$ such that
    \begin{equation}
        \widetilde{Y}_{i}(t_i)[q]\overset{d}{\rightarrow} N(0,1) \quad \quad \text{ as } t_i \to \infty,
    \end{equation}
    then
        \begin{equation}
         \widetilde{Y}(t_1,\ldots,t_p)[q]\overset{d}{\rightarrow} N(0,1) \quad \quad \text{ as } t_1,\ldots,t_p \to \infty.
        \end{equation}
\end{proposition}

    Let $B=(B_x)_{x\in\R^d}$ be a real-valued, continuous, centered, stationary Gaussian field with unit-variance. Let $C:\R^d \rightarrow \R$ be the covariance function of $B$ and assume it is separable in the sense of \eqref{equ:sep}. Then \begin{equation}
        \var(Y(t_1,\ldots,t_p)[q]) = (q!)^{1-p} \prod_{i=1}^p \var (Y_i{(t_i)}[q]).
    \end{equation}

    \begin{align} \label{equ:contrHq}
    & \|f{(t_1,\dots,t_p)} \otimes_r  f{(t_1,\dots,t_p)}\|_{{2q-2r}}^2 = \\
    & = \int_{(t_1D_1\times\dots\times t_pD_p)^4} C(x-z)^rC(y-u)^rC(x-y)^{q-r}C(z-u)^{q-r}dxdydudz  \\
    & =  \prod_{i=1}^p \int_{(t_iD_i)^4} C_i(x_i-z_i)^rC_i(y_i-u_i)^rC_i(x_i-y_i)^{q-r}C_i(z_i-u_i)^{q-r}dx_idy_idu_idz_i ,
\end{align}

where $x_i$ denotes the projection of $x$ onto the $i$th block of $\mathbb R^{d}=\prod_{i=1}^p\mathbb R^{d_i}$. Let us define $\tilde f:= \frac{f}{\|f\|}$. By linearity, it immediately follows that $\widetilde Y{(t_1,\dots,t_p)}[q]=I_q(\tilde f{(t_1,\dots,t_p)})$ and $\widetilde Y_i{(t_i)}[q]=I_q(\tilde f_i{(t_i)})$. Then, combining \eqref{equ:contrHq} and \eqref{equ:varHq}, we obtain \begin{equation}\label{equ:contr}
    \|\widetilde f{(t_1,\ldots,t_p)}\otimes_r \widetilde f{(t_1,\ldots,t_p)}\|_{2q-2r}^2 =  \prod_{i=1}^p  \|\widetilde f_i{(t_i)}\otimes_r \widetilde f_i(t_i)\|_{2q-2r}^2.
\end{equation}

Since for every $f\in \mathcal H^{\otimes q}$ and $r=1,\ldots,q-1$, we have \begin{equation}\label{equ:aiuto}
    \|f\otimes_r f\|_{2q-2r} ^2 \le \|f\|_{q}^4,
\end{equation}

\begin{proposition}\label{prop:quantHq}
    Let the same notations and assumptions of Proposition~\ref{prop:suffHq} prevail. In particular $N\sim N(0,1)$.
    Then, the following estimate holds: 
    \begin{align}\label{equ:quantHqLe}
        &d_{TV}(\widetilde Y{(t_1,\ldots,t_p)}[q],N)\le c_q \, \prod_{i=1}^p \sqrt{\E\big[\widetilde Y_i[q]^4\big]-3}, \\
        &{d_{W}(\widetilde Y{(t_1,\ldots,t_p)}[q],N)\le \frac{c_q}{\sqrt{\pi}} \, \prod_{i=1}^p \sqrt{\E\big[\widetilde Y_i[q]^4\big]-3},}
    \end{align}
    where $c_q=\sqrt{\frac{4}{q}\sum_{r=1}^{q-1}rr!^2\binom{q}{r}^4(2q-2r)!}$.
\end{proposition}

\begin{proof}[Proof of Proposition~\ref{prop:quantHq}]
    Recall the result of Theorem \ref{thm:4thQuant}, namely
    \begin{equation}\label{coco}
    d_{TV}(\widetilde Y{(t_1,\ldots,t_p)}[q],N)\le \frac{2}{\sqrt{3}} \sqrt{\E\big[(\widetilde Y{(t_1,\ldots,t_p)}[q])^4\big]-3},
    \end{equation}
    and
    \begin{align}
        & \E\big[(\widetilde Y{(t_1,\ldots,t_p)}[q])^4\big]-3 \label{coco2}\\
        & = \frac3q\sum_{r=1}^{q-1} rr!^2\binom{q}{r}^4(2q-2r)!
\|\widetilde f(t_1\ldots,t_p)\widetilde{\otimes}_r \widetilde f(t_1\ldots,t_p)\|^2_{2q-2r}\\
        &= \sum_{r=1}^{q-1} q!^2\binom{q}{r}^2\bigg\{\|\widetilde f(t_1\ldots,t_p)\otimes_r \widetilde f(t_1\ldots,t_p)\|^2_{2q-2r} \\
        & \hspace{6cm}
        +\binom{2q-2r}{q-r}\|\widetilde f(t_1\ldots,t_p)\widetilde{\otimes}_r \widetilde f(t_1\ldots,t_p)\|^2_{2q-2r}\bigg\}.
    \end{align}
    We deduce 
    \begin{align}
         d_{TV}(\widetilde Y{(t_1,\ldots,t_p)}[q],N)
        &\le \frac{2}{\sqrt{3}} \sqrt{\E\big[(\widetilde Y{(t_1,\ldots,t_p)}[q])^4\big]-3}\\
        &\le c_q\,\max_{r=1,\ldots,q-1} \|\widetilde f(t_1\ldots,t_p)\otimes_r  \widetilde  f(t_1\ldots,t_p)\|^2_{2q-2r}.
    \end{align}
    But, thanks to \eqref{equ:contr} and recalling \eqref{equ:aiuto}, 
   we have 
    \begin{align}
        \max_{r=1,\ldots,q-1} \|\widetilde f(t_1\ldots,t_p)\otimes_r  \widetilde  f(t_1\ldots,t_p)\|^2_{2q-2r}
        &= \max_{r=1,\ldots,q-1} \prod_{i=1}^p \|\widetilde f_i(t_i)\otimes_r \widetilde f_i(t_i)\|^2_{2q-2r}\\
        &\le  \prod_{i=1}^p \max_{r=1,\ldots,q-1} \|\widetilde f_i(t_i)\otimes_r \widetilde f_i(t_i)\|^2_{2q-2r}.
    \end{align}
    Also, as a consequence of the second equality in \eqref{coco2} (with $\widetilde{Y}_i$ instead of $\widetilde{Y}$), we have
    $$
    \|\widetilde f_i(t_i)\otimes_r \widetilde f_i(t_i)\|^2_{2q-2r} \leq \E\big[\widetilde Y_i[q]^4\big]-3,
    $$
    and the desired conclusions now easily follows.
\end{proof}

\end{document}
